\section{Trabalhos Relacionados}

Em geral outros já realizaram o trabalho de avaliar diferentes formas de minerar
padrões de uso, destacam-se trabalhos de revisão na área como
\cite{api_property_inference}
\cite{recommender_system_software_engineering_nearest_neighbours}, que comparam
várias formas de obter informações sobre desenvolvimento de software analisando
grandes quantidades de código.

Frequentemente na área são propostas de metodologia distintas, diferindo em:
\begin{itemize}
  \item fontes de dados distintas \cite{up_miner} %TODO: add ref to mining of
                                                  % stack overflow data
  \item utilizar-se de aspectos da linguagem que analisam \cite{mapo}
  \item utilizar-se de informações do código em execução \cite{groums}
  \cite{focus}
\end{itemize}

Porém, aparte alguns casos especiais \cite{probabilistic_api_miner}
\cite{groums} \cite{prefix_span} \cite{focus}, o método mais comum de encontrar
padrões é através de mineração de regras de associação.

Um trabalho extremamente similar ao presente é o feito por H. Kagdi et al.
\cite{comp_mining_source_code}. O estudo avalia mineração de padrões de uso em
C, comparando mineração por regras de associação e por mineração de padrões
sequenciais. Além da escolha de métodos sendo comparados, tal análise minera e
avalia padrões de uso de todas funções no kernel linux. O número de violações
dos padrões de uso é utilizado como ponto de comparação para os resultados 

Adicionalmente, tal estudo avalia padrões de uso ordenados e não ordenados,
enquanto o presente avalia apenas padrões de uso não ordenados.

Como levantado em \cite{api_property_inference}
\cite{analyzing_temporal_api_usage_patterns}, a avaliação automática de
resultados na área é um problema recorrente. A decisão escolhida é
frequentemente apelar para simples análises descritivas dos resultados ou, se
houver tempo e orçamento o suficientes disponíveis, realizar entrevistas com
desenvolvedores para validar os dados encontrados. 
